\chapter{Selected Source Listings}\label{ch:ap-d}

This appendix reproduces the handful of Go functions that carry the most conceptual weight in the APEX backend. Each is quoted faithfully from the real source tree under \texttt{internal/}; the only edits are the occasional trimming of a long sibling function so the reader's attention stays on the mechanism under discussion. The listings are the concrete counterparts of the pipelines drawn abstractly in \cref{fig:grading-seq}, \cref{fig:authz}, and \cref{fig:reminder}. Reading them together should make it clear how little exotic machinery the system relies on: the correctness of APEX rests on a small number of carefully ordered database writes and a few honest normalisation and aggregation rules, not on a framework doing the thinking for us.

\section{The Grading Engine}\label{sec:ap-d-grader}

The grading engine in \texttt{internal/judge0/grader.go} is the heart of the project. It is deliberately plain Go: it loads rows, calls out to Judge0 over HTTP, compares strings, and writes results back. Three parts are worth reading closely --- the output-normalisation helper, the per-test-case grading loop, and the attempt-level aggregator.

\subsection{Output Normalisation}\label{sec:ap-d-normalize}

Comparing program output to a teacher's expected answer is where naive auto-graders quietly fail thousands of students. Judge0 returns whatever the sandboxed process wrote to \texttt{stdout}, which on many toolchains carries \texttt{\textbackslash r\textbackslash n} line endings and a trailing newline, while the expected output was typed into a web form and may carry stray trailing spaces. \texttt{normalizeOutput} in \cref{lst:normalize-app} applies exactly four transformations --- CRLF to LF, bare CR to LF, right-trim of spaces and tabs on every line, then a whole-string trim --- and is applied \emph{symmetrically} to both sides before the equality check. What to notice is the restraint: it does not collapse interior whitespace or lowercase anything, so a genuinely wrong answer still fails. It only forgives the differences that are artefacts of transport and data entry, never differences of substance.

\begin{lstlisting}[language=Go,caption={Symmetric output normalisation applied to both actual and expected output before comparison.},label={lst:normalize-app}]
// normalizeOutput strips CR, trailing whitespace on each line, and trims edges
// so Judge0 output (often \r\n with trailing newline) compares cleanly to
// teacher-entered expected output.
func normalizeOutput(s string) string {
	s = strings.ReplaceAll(s, "\r\n", "\n")
	s = strings.ReplaceAll(s, "\r", "\n")
	lines := strings.Split(s, "\n")
	for i, ln := range lines {
		lines[i] = strings.TrimRight(ln, " \t")
	}
	return strings.TrimSpace(strings.Join(lines, "\n"))
}
\end{lstlisting}

\subsection{The Coding Grading Loop}\label{sec:ap-d-grade}

\cref{lst:grade} is the core coding path, \texttt{Grade}. Several design decisions are visible in it. First, the empty-test-case guard: a coding problem with no configured test cases is scored \texttt{100.0} and marked \texttt{accepted} rather than crashing or scoring zero --- a defensive choice that avoids penalising students for a teacher's incomplete setup. Second, the status mapping treats Judge0 status \texttt{3} (Accepted) and \texttt{4} (Wrong Answer) identically as ``the program ran, now compare the strings''; only a genuine execution failure (\texttt{6} compilation error, \texttt{5} time-limit, \texttt{11} runtime error) short-circuits the comparison. Third, one \texttt{TestResult} row is persisted per test case with its own timing and memory figures, while the maximum time and memory across all tests are rolled up onto the \texttt{Submission}. Finally, the score is the simple fraction \texttt{passed / total * 100}, and control passes to \texttt{maybeFinalizeAttempt} to decide whether the whole exam attempt can now be finalised. The notification is only fired for standalone submissions (\texttt{ExamAttemptID == nil}); exam submissions defer their single notification to the aggregator.

\begin{lstlisting}[language=Go,caption={The coding grading loop: run each test case through Judge0, compare normalised output, persist per-test results, and roll up the score.},label={lst:grade}]
func Grade(submissionID uint) {
	var submission models.Submission
	if err := database.DB.First(&submission, submissionID).Error; err != nil {
		log.Printf("grading: submission %d not found: %v", submissionID, err)
		return
	}

	database.DB.Model(&submission).Update("status", "running")

	var testCases []models.TestCase
	database.DB.Where("problem_id = ?", submission.ProblemID).Order("order_index asc").Find(&testCases)

	if len(testCases) == 0 {
		database.DB.Model(&submission).Updates(map[string]any{
			"status": "accepted",
			"score":  100.0,
		})
		return
	}

	passed := 0
	total := len(testCases)
	var maxTimeMs int
	var maxMemoryKb int

	for _, tc := range testCases {
		result := RunCode(submission.Code, submission.Language, tc.Input)

		actualOutput := normalizeOutput(result.Stdout)
		expectedOutput := normalizeOutput(tc.ExpectedOutput)
		// StatusID 3 = Accepted; treat StatusID 4 (Wrong Answer) as valid-run-but-compare.
		ran := result.StatusID == 3 || result.StatusID == 4
		isPassed := ran && actualOutput == expectedOutput

		if isPassed {
			passed++
		}

		status := "accepted"
		if !isPassed {
			switch result.StatusID {
			case 6:
				status = "compilation_error"
			case 5:
				status = "time_limit_exceeded"
			case 11:
				status = "runtime_error"
			default:
				status = "wrong_answer"
			}
		}

		tr := models.TestResult{
			SubmissionID:    submission.ID,
			TestCaseID:      tc.ID,
			Passed:          isPassed,
			ActualOutput:    actualOutput,
			ExecutionTimeMs: result.TimeMs,
			MemoryKb:        result.MemoryKb,
			Status:          status,
		}
		database.DB.Create(&tr)

		if result.TimeMs > maxTimeMs {
			maxTimeMs = result.TimeMs
		}
		if result.MemoryKb > maxMemoryKb {
			maxMemoryKb = result.MemoryKb
		}
	}

	score := float64(passed) / float64(total) * 100.0
	finalStatus := "wrong_answer"
	if passed == total {
		finalStatus = "accepted"
	}

	database.DB.Model(&submission).Updates(map[string]any{
		"status":            finalStatus,
		"passed_count":      passed,
		"total_count":       total,
		"score":             score,
		"execution_time_ms": maxTimeMs,
		"memory_kb":         maxMemoryKb,
	})

	if submission.ExamAttemptID == nil {
		notification.Create(submission.UserID, "result",
			"Submission Graded",
			fmt.Sprintf("Your coding submission scored %.0f%% (%d/%d test cases passed)", score, passed, total),
			"/dashboard/results",
		)
	}
	maybeFinalizeAttempt(submission.ExamAttemptID)
}
\end{lstlisting}

\subsection{Attempt Aggregation}\label{sec:ap-d-finalize}

A single exam attempt spans many submissions of mixed type (coding, MCQ, written), each graded independently in its own goroutine. \texttt{maybeFinalizeAttempt} in \cref{lst:finalize} is the join point that turns those scattered results into one attempt score, and it is written to be safe to call after \emph{every} submission finishes. Three subtleties deserve attention. First, it returns early if any sibling submission is still \texttt{pending} or \texttt{running}, so the last grader to finish is the one that actually aggregates. Second, it iterates the exam's \emph{full problem list} rather than the submissions, so a question the student skipped correctly contributes its points to the denominator at a score of zero. Third, submissions in \texttt{pending\_review} --- written answers awaiting manual grading --- are excluded from both numerator and denominator, so an ungraded essay does not artificially depress the visible total before the teacher gets to it. The ``Exam Graded'' notification is fired exactly once, guarded by the \texttt{graded\_notified} column, and only after no \texttt{pending\_review} submissions remain.

\begin{lstlisting}[language=Go,caption={Attempt-level aggregation: re-score over the exam's problem list once no submission is still in flight, excluding pending manual grades.},label={lst:finalize}]
// maybeFinalizeAttempt aggregates score across all submissions once none
// remain in pending/running state. Safe to call multiple times.
func maybeFinalizeAttempt(attemptID *uint) {
	if attemptID == nil {
		return
	}
	var pending int64
	database.DB.Model(&models.Submission{}).
		Where("exam_attempt_id = ? AND status IN ?", *attemptID, []string{"pending", "running"}).
		Count(&pending)
	if pending > 0 {
		return
	}

	var attempt models.ExamAttempt
	if err := database.DB.First(&attempt, *attemptID).Error; err != nil {
		return
	}

	var problems []models.Problem
	database.DB.Where("exam_id = ?", attempt.ExamID).Find(&problems)
	if len(problems) == 0 {
		return
	}

	var subs []models.Submission
	database.DB.Where("exam_attempt_id = ?", *attemptID).Find(&subs)
	subByProblem := make(map[uint]models.Submission, len(subs))
	for _, s := range subs {
		subByProblem[s.ProblemID] = s
	}

	// Iterate the exam's problem list so skipped questions (no submission)
	// stay in the denominator at score=0. Pending-review submissions
	// (written answers awaiting manual grading) are still excluded from
	// both numerator and denominator until a teacher grades them.
	var earnedPoints float64
	var totalPoints float64
	for _, p := range problems {
		s, hasSub := subByProblem[p.ID]
		if hasSub && s.Status == "pending_review" {
			continue
		}
		pts := float64(p.Points)
		if pts <= 0 {
			pts = 10
		}
		totalPoints += pts
		if hasSub {
			earnedPoints += s.Score / 100.0 * pts
		}
	}
	pct := 0.0
	if totalPoints > 0 {
		pct = earnedPoints / totalPoints * 100.0
	}

	database.DB.Model(&models.ExamAttempt{}).
		Where("id = ?", *attemptID).
		Update("score", pct)

	// Fire a single "Exam Graded" notification once all submissions are
	// fully graded (no pending_review remaining). Idempotent via
	// exam_attempts.graded_notified.
	var anyPendingReview int64
	database.DB.Model(&models.Submission{}).
		Where("exam_attempt_id = ? AND status = ?", *attemptID, "pending_review").
		Count(&anyPendingReview)
	if anyPendingReview > 0 {
		return
	}

	if err := database.DB.Preload("Exam").First(&attempt, *attemptID).Error; err != nil {
		return
	}
	if attempt.Status != "submitted" || attempt.GradedNotified {
		return
	}
	database.DB.Model(&attempt).Update("graded_notified", true)
	notification.Create(attempt.UserID, "result",
		"Exam Graded",
		fmt.Sprintf("Your submission for \"%s\" has been graded. Score: %.0f%%", attempt.Exam.Title, pct),
		fmt.Sprintf("/dashboard/exam/%d/review", attempt.ExamID),
	)
}
\end{lstlisting}

\section{Authentication and Authorisation Middleware}\label{sec:ap-d-auth}

The first two of APEX's three authorisation gates (\cref{fig:authz}) live in \texttt{internal/middleware/auth.go}, shown in \cref{lst:auth}. \texttt{Auth} validates the bearer token and is defensive at every step that involves untrusted input. The key-lookup callback rejects any token whose signing method is not \texttt{*jwt.SigningMethodHMAC}, closing the classic ``\texttt{alg:none}'' and algorithm-confusion attacks in which an attacker asks the server to verify an HS256 token using an RS256 public key as the HMAC secret \cite{owasp-jwt}. Every claim is then read with a comma-ok type assertion --- \texttt{user\_id} as \texttt{float64} (JSON's only number type, cast back to \texttt{uint}), \texttt{email} and \texttt{role} as \texttt{string} --- so a malformed or partial token produces a clean 401 rather than a panic. Only after all checks pass are the three values written onto the Gin context for downstream handlers to trust. \texttt{RequireRole} is the second gate: it reads the \texttt{role} the first gate deposited, returns 401 if it is absent (an unauthenticated request), and 403 ``insufficient permissions'' if it is present but not among the allowed roles \cite{rfc7519,rfc6749}.

\begin{lstlisting}[language=Go,caption={JWT authentication middleware and role gate: signing-method enforcement, safe claim extraction, and the 401-vs-403 distinction.},label={lst:auth}]
func Auth(secret string) gin.HandlerFunc {
	return func(c *gin.Context) {
		header := c.GetHeader("Authorization")
		if header == "" || !strings.HasPrefix(header, "Bearer ") {
			c.AbortWithStatusJSON(http.StatusUnauthorized, gin.H{"error": "missing or invalid token"})
			return
		}

		tokenStr := strings.TrimPrefix(header, "Bearer ")
		token, err := jwt.Parse(tokenStr, func(t *jwt.Token) (any, error) {
			if _, ok := t.Method.(*jwt.SigningMethodHMAC); !ok {
				return nil, jwt.ErrSignatureInvalid
			}
			return []byte(secret), nil
		})
		if err != nil || !token.Valid {
			c.AbortWithStatusJSON(http.StatusUnauthorized, gin.H{"error": "invalid token"})
			return
		}

		claims, ok := token.Claims.(jwt.MapClaims)
		if !ok {
			c.AbortWithStatusJSON(http.StatusUnauthorized, gin.H{"error": "invalid token claims"})
			return
		}

		// Safe type assertions with comma-ok pattern.
		userIDFloat, ok := claims["user_id"].(float64)
		if !ok {
			c.AbortWithStatusJSON(http.StatusUnauthorized, gin.H{"error": "invalid user_id in token"})
			return
		}
		email, ok := claims["email"].(string)
		if !ok {
			c.AbortWithStatusJSON(http.StatusUnauthorized, gin.H{"error": "invalid email in token"})
			return
		}
		role, ok := claims["role"].(string)
		if !ok {
			c.AbortWithStatusJSON(http.StatusUnauthorized, gin.H{"error": "invalid role in token"})
			return
		}

		c.Set("user_id", uint(userIDFloat))
		c.Set("email", email)
		c.Set("role", role)
		c.Next()
	}
}

func RequireRole(roles ...string) gin.HandlerFunc {
	return func(c *gin.Context) {
		role, exists := c.Get("role")
		if !exists {
			c.AbortWithStatusJSON(http.StatusUnauthorized, gin.H{"error": "unauthorized"})
			return
		}
		roleStr, ok := role.(string)
		if !ok {
			c.AbortWithStatusJSON(http.StatusUnauthorized, gin.H{"error": "invalid role"})
			return
		}
		for _, r := range roles {
			if roleStr == r {
				c.Next()
				return
			}
		}
		c.AbortWithStatusJSON(http.StatusForbidden, gin.H{"error": "insufficient permissions"})
	}
}
\end{lstlisting}

The third gate --- resource ownership --- deliberately does not live in middleware, because whether a user ``owns'' a resource depends entirely on the resource. It is enforced inside each handler against the \texttt{user\_id} the middleware deposited, as \cref{fig:authz} illustrates.

\section{The Reminder Scheduler}\label{sec:ap-d-reminder}

APEX runs its exam reminders in-process rather than depending on an external cron system, keeping the deployment to a single Go binary. \cref{lst:scheduler} shows the driver in \texttt{internal/reminder/scheduler.go}. \texttt{Start} spawns one goroutine that waits fifteen seconds for the database connection pool to settle, then ticks once a minute forever. The \texttt{runOnce} tick is wrapped in a deferred \texttt{recover} so that a panic in one sweep --- a malformed row, a transient query error --- logs and is swallowed rather than killing the scheduler goroutine and silently stopping all future reminders. Each tick performs the in-app T-60 sweep inline and then delegates the two email sweeps. The critical property, visible in the \texttt{WHERE} clauses of \cref{fig:reminder}, is that every sweep filters on a dedicated ``sent-at'' timestamp column being \texttt{NULL} and stamps it after sending. Because the dedup state lives in the database rather than in memory, a process restart cannot cause a reminder to be re-sent: the scheduler is idempotent across restarts.

\begin{lstlisting}[language=Go,caption={The reminder scheduler tick loop: a fifteen-second warm-up, a one-minute ticker, and a panic-guarded sweep over published exams.},label={lst:scheduler}]
// Start launches a goroutine that periodically scans for upcoming exams and
// notifies enrolled students who opted in to exam reminders. Reminder fires
// once per exam (deduped via exams.reminder_sent_at).
func Start() {
	go func() {
		// Initial short delay to let DB settle.
		time.Sleep(15 * time.Second)
		tick := time.NewTicker(1 * time.Minute)
		defer tick.Stop()
		runOnce()
		for range tick.C {
			runOnce()
		}
	}()
}

func runOnce() {
	defer func() {
		if r := recover(); r != nil {
			log.Printf("[reminder] panic: %v", r)
		}
	}()

	now := time.Now()
	cutoff := now.Add(60 * time.Minute)

	var exams []models.Exam
	if err := database.DB.
		Where("is_draft = ? AND reminder_sent_at IS NULL AND start_time IS NOT NULL AND start_time > ? AND start_time <= ?", false, now, cutoff).
		Find(&exams).Error; err != nil {
		log.Printf("[reminder] query failed: %v", err)
		return
	}

	for _, exam := range exams {
		notifyExam(exam)
	}

	emailSweep1Hour(now)
	emailSweepStart(now)
}
\end{lstlisting}

The fan-out itself, in \texttt{notifyExam}, resolves an exam to its assigned classes, those classes to their enrolled members, and those members to the subset whose profile has \texttt{notify\_exam\_reminders = true}. One nuance worth noting is that an exam with no assigned classes or no enrolled members is still stamped \texttt{reminder\_sent\_at} so that it is not rescanned on every subsequent tick for the rest of its life --- a small guard that keeps the sweep's working set bounded to genuinely pending exams.

\section{Destructive-Edit Reset}\label{sec:ap-d-reset}

When a teacher makes a destructive edit to a live exam --- changing a test case, editing an MCQ's correct answers, or deleting a problem --- every attempt already scored against the old configuration is invalid. \texttt{ResetExamAttempts} in \texttt{internal/exam/reset.go}, shown in \cref{lst:reset}, cleans this up atomically. The whole operation runs inside a single \texttt{database.DB.Transaction}, so either all of the student progress is wiped and the exam re-stamped, or nothing changes. The deletion order respects foreign-key dependencies: \texttt{TestResult} rows are removed first (looked up via their parent submission IDs), then \texttt{Submission} rows, then the \texttt{ExamAttempt} rows themselves. The transaction then stamps \texttt{exams.reset\_at}, un-announces grades at the class level so stale scores stop showing, and --- only if there were real attempts to reset --- notifies enrolled students that they may re-enter. That last guard is what keeps the ``replace-all problems'' save path from fanning out one spurious notification per deleted problem.

\begin{lstlisting}[language=Go,caption={Transactional reset of student progress after a destructive teacher edit, with FK-ordered deletes and a guarded notification.},label={lst:reset}]
func ResetExamAttempts(examID uint) error {
	return database.DB.Transaction(func(tx *gorm.DB) error {
		var exam models.Exam
		if err := tx.First(&exam, examID).Error; err != nil {
			return err
		}

		var attemptCount int64
		tx.Model(&models.ExamAttempt{}).Where("exam_id = ?", examID).Count(&attemptCount)
		hadAttempts := attemptCount > 0

		var subIDs []uint
		tx.Model(&models.Submission{}).Where("exam_id = ?", examID).Pluck("id", &subIDs)
		if len(subIDs) > 0 {
			if err := tx.Where("submission_id IN ?", subIDs).Delete(&models.TestResult{}).Error; err != nil {
				return err
			}
		}
		if err := tx.Where("exam_id = ?", examID).Delete(&models.Submission{}).Error; err != nil {
			return err
		}
		if err := tx.Where("exam_id = ?", examID).Delete(&models.ExamAttempt{}).Error; err != nil {
			return err
		}

		now := time.Now()
		if err := tx.Model(&models.Exam{}).Where("id = ?", examID).Update("reset_at", now).Error; err != nil {
			return err
		}
		exam.ResetAt = &now

		// Un-announce grades for the classes this exam belongs to — old
		// announced scores are now invalid.
		var classIDs []uint
		tx.Model(&models.ExamClass{}).Where("exam_id = ?", examID).Pluck("class_id", &classIDs)
		if len(classIDs) > 0 {
			if err := tx.Model(&models.Class{}).Where("id IN ?", classIDs).Update("grades_announced", false).Error; err != nil {
				return err
			}
		}

		// Only notify if there were actual attempts to reset.
		if !hadAttempts {
			return nil
		}
		if len(classIDs) > 0 {
			var userIDs []uint
			tx.Model(&models.ClassMember{}).Where("class_id IN ?", classIDs).Pluck("user_id", &userIDs)
			link := fmt.Sprintf("/dashboard/exam/%d", examID)
			desc := fmt.Sprintf("Exam '%s' was updated by your teacher. Your attempt has been reset — you can re-enter.", exam.Title)
			for _, uid := range userIDs {
				notification.Create(uid, "exam_reset", "Exam updated", desc, link)
			}
		}
		return nil
	})
}
\end{lstlisting}

Taken together, these five listings capture the engineering personality of APEX: correctness enforced by explicit, ordered database writes; safety enforced by early returns and comma-ok assertions rather than exceptions; and idempotence enforced by persisted state columns rather than in-memory flags. There is no hidden magic --- which is precisely what makes the system auditable.
